\noindent On the fresh synthetic panels, the lifted output layer helps the transformer (Table~\ref{tab:v3}). With two
missing sensors the LCT beats the transformer with the identical backbone, recipe and seed by
\VthreeGOnekTwolctvspfnabs\ \VthreeGOnekTwolctvspfnci\ nats on G1 and \VthreeGThreekTwolctvspfnabs\
\VthreeGThreekTwolctvspfnci\ on sixteen sensors (E5, E6 pass). It matches the small lifted network on five sensors
(\VthreeGOnekTwolctvslift\ \VthreeGOnekTwolctvsliftci) and exceeds it on sixteen (\VthreeGThreekTwolctvslift\
\VthreeGThreekTwolctvsliftci), and it learns faster: its training NLL reaches the transformer's final value after
\Trainlctreach\ of 30{,}000 steps (Figure~\ref{fig:training}). On the v2 panels (descriptive) it closes \Gaplct\% of
the F1 gap (lifted network \GapliftOne\%) but inherits part of the transformer's degradation with 16 calibration rows.
After fine-tuning on Beijing NO$_2$, it beats the fine-tuned lifted network by \VthreebeijingnoTwoeZerolctvslift\
\VthreebeijingnoTwoeZerolctvsliftci\ nats (E7a passes), recovering \VthreebeijingnoTwoeZerorecovered\% of the gap
between the lifted network (\VthreebeijingnoTwoeZerovallift) and the transformer (\VthreebeijingnoTwoeZerovalpfn).
Against the transformer the difference is \VthreebeijingnoTwoeZerolctvspfn\ \VthreebeijingnoTwoeZerolctvspfnci; the
interval includes zero but extends below the margin of $-0.02$, so non-inferiority (E7b) is not established. With six
more stations missing, the two tie (\VthreebeijingnoTwoeSixlctvspfn\ \VthreebeijingnoTwoeSixlctvspfnci).
On Beijing CO (descriptive) the pattern repeats: the LCT (\VthreebeijingcoeZerovallct) recovers
\VthreebeijingcoeZerorecovered\% of the gap between the lifted network (\VthreebeijingcoeZerovallift) and the
transformer (\VthreebeijingcoeZerovalpfn).


